% from ega4-2-fr.tex
\oldpage[IV-2]{123}

il existe un entier $r$ tel que $\mathfrak m^rN=0$; si on pose
$N'_j=\mathfrak m^jN$ pour $0\leq j\leq r$, $N'_j/N'_{j+1}$ est un
$(B/\mathfrak m)$-module de type fini, donc somme directe d'un nombre fini de
sous-modules isomorphes à $B/\mathfrak m$, puisque $B/\mathfrak m$ est un
corps; autrement dit, il y a une filtration finie $(N_h)_{0\leq h\leq m}$ de
$N$ telle que $N_m=0$ et que $N_h/N_{h+1}$ soit isomorphe à
$B/\mathfrak m$, i.e. au corps des fractions de $A/\mathfrak p$; la
filtration des $M_h=M\cap N_h$ répond à la question, car $M_h/M_{h+1}$ est
isomorphe à un sous-$(A/\mathfrak p)$-module de type fini de
$N_h/N_{h+1}=B/\mathfrak m$; mais on sait qu'un tel sous-module est isomorphe
à un sous-module de $A/\mathfrak p$.

L'existence de la filtration $(M_h)$ montre alors, par le même raisonnement que
plus haut, que pour prouver (moyennant b)) que
$\mathcal H^0_{X/Z}(\mathcal F)$ est cohérent, on peut se borner à montrer que
$\mathcal H^0_{X/Z}(\mathcal G)$ est cohérent, où
$\mathcal G=(A/\mathfrak p)^\sim$, $\mathfrak p$ étant un idéal associé à
$M$. Mais si $\mathfrak p=\mathfrak j_y$ avec $y\in Z$, le support de
$\mathcal G$ est un ensemble fermé contenu dans $Z$, puisque $Z$ est stable
par spécialisation; la définition (5.9.1.2) montre alors que
$\mathcal H^0_{X/Z}(\mathcal G)=0$. Si au contraire
$\mathfrak p=\mathfrak j_{x_\alpha}$ pour un $\alpha$, on a
$\mathcal G=(h_\alpha)_*(\mathcal O_{Y_\alpha})$ par définition, et l'on a vu
plus haut que $\mathcal H^0_{X/Z}(\mathcal G)$ est isomorphe à
$(h_\alpha)_*(\mathcal H^0_{Y_\alpha/Z_\alpha}(\mathcal O_{Y_\alpha}))$,
donc est cohérent en vertu de l'hypothèse b).

On a déjà vu (5.10.10, (ii)) que a) entraîne c). Il reste à prouver que, sous
l'une ou l'autre des hypothèses (i), (ii), c) entraîne b). Notons que si $X$
vérifie (i), il en est de même de chaque $Y_\alpha$. Il suffit donc (5.9.3.1)
de prouver le

\begin{corollary}[5.11.2]
\EGAReaderTarget{ega4.c04.cor.5-11-2}{IV.2.s5.11.2-fr}
Soient $A$ un anneau noethérien intègre, vérifiant l'une des deux hypothèses
suivantes:
\begin{enumerate}
\item[(i)] $A$ est quotient d'un anneau noethérien régulier.
\item[(ii)] $A$ est universellement caténaire, et sa clôture intégrale $A'$ est
une $A$-algèbre finie.
\end{enumerate}
Alors l'anneau $A^{(1)}$, intersection des $A_\mathfrak p$ où $\mathfrak p$
parcourt l'ensemble des idéaux premiers de $A$ de hauteur 1, est une
$A$-algèbre finie.
\end{corollary}

(i) Posons $X=\operatorname{Spec}(A)$. L'ensemble $U$ des points $x\in X$ où
$X$ vérifie $(S_2)$ est ouvert dans $X$ moyennant l'hypothèse (i) (6.11.2)
\footnote{Le lecteur pourra vérifier que (5.11.2) n'est pas utilisé dans la
démonstration de (6.11.2).}. En outre, si l'on pose $Z=X-U$, on a
$\operatorname{codim}(Z,X)\geq2$: en effet, pour tout $x\in X$ tel que
$\dim(A_x)\leq1$, on a $\dim(A_{x'})\leq1$ pour tout $x'$ générisation de
$x$, et comme $A_{x'}$ est intègre, $\operatorname{prof}(A_{x'})\geq1$, donc
$X$ vérifie $(S_2)$ au point $x$. On a donc $Z\subset Z^{(2)}$, et
$\mathcal H^0_{X/Z^{(2)}}(\mathcal O_X)
=j_*(\mathcal H^0_{U/Z^{(2)}}(\mathcal O_U))$, où $j:U\to X$ est l'injection
canonique (5.9.1.2). Mais comme le préschéma $U$ vérifie $(S_2)$, il résulte de
(5.10.14) que $\mathcal H^0_{U/Z^{(2)}}(\mathcal O_U)$ est isomorphe à
$\mathcal O_U$; d'autre part, comme $\operatorname{codim}(Z,X)\geq2$, on sait,
d'après le chapitre III, \S~9, que $j_*(\mathcal O_U)$ est un
$\mathcal O_X$-Module cohérent, ce qui démontre la proposition dans le cas
(i).

(ii) L'anneau $A'$ est, en vertu de l'hypothèse (ii), un anneau noethérien
intègre et intégralement clos, donc (Bourbaki, \emph{Alg. comm.}, chap. VII,
\S~1, n\textsuperscript{o}~6, th. 4) intersection de ses anneaux locaux
$A'_{\mathfrak p'}$, où $\mathfrak p'$ parcourt l'ensemble des idéaux premiers
de hauteur 1 de $A'$. Or, pour un tel idéal premier $\mathfrak p'$, si l'on
pose $\mathfrak p=\mathfrak p'\cap A$ et $S=A-\mathfrak p$,
$A'_{\mathfrak p'}$ est un anneau local en l'idéal premier
$S^{-1}\mathfrak p'$ de $S^{-1}A'$, et $S^{-1}A'$ est par hypothèse

